\section{Theoretical Analysis}
\label{sec:theory}

\subsection{Main Regret Bound}

\begin{theorem}[Regret bound for DVA-MCTS]
\label{thm:main}
Under Assumptions~\ref{ass:lipschitz}--\ref{ass:noise}, with threshold
$\tau_t = \gamma / \sqrt{t}$ and dynamic Lipschitz proxy
(Definition~\ref{def:dynamic_proxy}), DVA-MCTS achieves:
\begin{enumerate}[(i)]
  \item \textbf{(Single-level tree, $D{=}1$).}
        \[
          \mathbb{E}[R(T)] \;\leq\; C\sqrt{T\log K} \;+\; 2\gamma\sqrt{T}
                                   \;+\; LN_{\mathcal{T}},
          \qquad C = 8\sigma_{\max}\sqrt{2},
        \]
        with $\gamma^* = 4\sigma_{\max}\sqrt{2\log K}$ giving
        $\mathbb{E}[R(T)] \leq 16\sigma_{\max}\sqrt{2}\sqrt{T\log K} + LN_{\mathcal{T}}$.
        This matches the $\Omega(\sqrt{T})$ lower bound (Theorem~\ref{thm:lower})
        up to the $\sqrt{\log K}$ factor and the constant $LN_{\mathcal{T}}$.
  \item \textbf{(Depth-$D$ tree, general $D \geq 1$).}
        \[
          \mathbb{E}[R(T)] \;\leq\; C(K,D)\sqrt{T\log K} \;+\; 2\gamma\sqrt{T}
                                   \;+\; LDN_{\mathcal{T}},
        \]
        where $C(K,D) = 16\sigma_{\max}\sqrt{2}\,\Phi(K,D)$ and
        $\Phi(K,D) = (K^{D/2}-1)/(\sqrt{K}-1)$ is the tree-geometry constant
        from Lemma~\ref{lem:uct}.  For fixed $K, D$: $LDN_{\mathcal{T}}$ is a constant
        independent of $T$, so the rate is $O(\sqrt{T\log K})$;
        for $K{=}2$, $D{=}12$: $\Phi(2,12) \leq 153$.
\end{enumerate}
The $O(\sqrt{T\log K})$ rate holds in both parts.
Part (ii) inherits a depth-dependent constant $\Phi(K,D)$ from the tree geometry
(Lemma~\ref{lem:uct}; see Remark~\ref{rem:tree_depth}) and an additive
$LDN_{\mathcal{T}}$ constant from the dynamic proxy estimation error
(Remark~\ref{rem:LDN_T}).  Both are $O(1)$ for fixed $K, D$.
\end{theorem}

\begin{definition}[Dynamic Lipschitz proxy]
\label{def:dynamic_proxy}
At each step $t$, for any node $s$, the \emph{dynamic proxy estimate} is
\[
  \hat{v}_t(s) \;=\;
  \begin{cases}
    \hat{V}(s) & \text{if $s$ has been verified by step $t$,}\\
    \hat{V}(s_{\mathrm{anc},t}(s)) & \text{otherwise,}
  \end{cases}
\]
where $s_{\mathrm{anc},t}(s)$ is the deepest ancestor of $s$ verified by step $t$.
The proxy is recomputed fresh at each step---not cached---so the estimation error
$|V(s) - \hat{v}_t(s)| \leq L\cdot\Delta d_t(s) + \sigma_{\max}$ reflects the
\emph{current} tree coverage, improving automatically as ancestors are verified.
\end{definition}

\begin{remark}[Tightness of the constant]
\label{rem:constant_tightness}
The constant $C = 16\sigma_{\max}\sqrt{2}$ in part~(i) is \emph{order-optimal}
relative to the lower bound: Theorem~\ref{thm:lower} gives
$\mathbb{E}[R(T)] \geq (\sigma/8)\sqrt{T}$, and the ratio is
$C/(\sigma/8) = 128\sqrt{2} \approx 181$, independent of $T$ or $K$.
This gap arises from the union-bound and Cauchy--Schwarz steps that are tight
only in adversarial worst-case configurations where all nodes are visited equally.
UCT concentrates exploration on high-value paths, so the empirical constant is
$3$--$10\times$ smaller (Figure~\ref{fig:regret}(b)).
For depth $D > 1$, the factor $\Phi(K,D)$ is a Cauchy--Schwarz artifact.
A matching lower bound for general $D$ incorporating tree geometry is an open
problem (see Remark~\ref{rem:lb_scope} in Appendix~\ref{app:lower}).

The additional additive term $LDN_{\mathcal{T}}$ in Theorem~\ref{thm:main}(ii) is
a constant (for fixed $K, D$) that arises from the dynamic proxy estimation error
accumulated while oracle-best nodes are unverified.  It does not affect the
asymptotic rate but may dominate for small $T$ (see Remark~\ref{rem:LDN_T}).
For the small-tree setting ($K{=}2$, $D{=}4$, $N_{\mathcal{T}}{=}30$),
$LDN_{\mathcal{T}} = 0.35 \times 4 \times 30 = 42$, negligible relative to observed
empirical regret at all tested budgets.
\end{remark}

\begin{proof}[Proof sketch]
We decompose the regret at each step $t$ into two parts:
\[
  V(s^*_t) - V(\hat{s}^*_t)
  \;=\;
  \underbrace{V(s^*_t) - \hat{v}_t(s^*_t)}_{\text{(I) estimation error}}
  \;+\;
  \underbrace{\hat{v}_t(s^*_t) - V(\hat{s}^*_t)}_{\text{(II) selection error}}.
\]
Here $\hat{v}_t(s)$ is the \emph{dynamic} proxy estimate (Definition~\ref{def:dynamic_proxy}),
recomputed at each step $t$ using the most recently verified ancestor of $s$.

\textbf{Bounding term (I): estimation error for $s^*_t$.}
Since $s^*_t$ is a visited node, its dynamic estimate $\hat{v}_t(s^*_t)$ is either its
own verified score or its nearest verified ancestor's score.

\emph{Case A ($s^*_t$ verified by step $t$):}
$\mathbb{E}[V(s^*_t) - \hat{v}_t(s^*_t)] = \mathbb{E}[-\varepsilon_{s^*_t}] = 0$.

\emph{Case B ($s^*_t$ not yet verified by step $t$):}
$\mathbb{E}[V(s^*_t) - \hat{v}_t(s^*_t)] \leq L\cdot\Delta d_t(s^*_t) \leq LD$.
By the single-verification property (Lemma~\ref{lem:single_verify}),
$s^*_t$ can remain unverified for at most $N_{\mathcal{T}}$ of the $T$ steps.
Hence the total Case-B contribution is $\leq LDN_{\mathcal{T}}$, a constant.

\textbf{Bounding term (II): selection error.}
Since $\hat{s}^*_t = \arg\max_{i \leq t}\hat{v}_t(s_i)$, we have
$\hat{v}_t(s^*_t) \leq \hat{v}_t(\hat{s}^*_t)$, so term~(II) is the overestimation
at $\hat{s}^*_t$.  Lemma~\ref{lem:uct} (backward induction over $D$ tree levels,
Appendix~\ref{app:proofs}) gives cumulative selection error
$\leq C_1(K,D)\sqrt{T\log K}$.

\textbf{Combining.}
\[
  \mathbb{E}[R(T)]
  \;\leq\; C_1(K,D)\sqrt{T\log K} + LDN_{\mathcal{T}}.
\]
The threshold schedule additionally introduces a $2\gamma\sqrt{T}$ term from
current-step proxy error (Lemma~\ref{lem:proxy}); setting $\gamma^*$ to balance
terms gives $C(K,D) = 2C_1(K,D) = 16\sigma_{\max}\sqrt{2}\,\Phi(K,D)$.
The full proof is in Appendix~\ref{app:proofs}.
\end{proof}

\subsection{Verifier Call Complexity}

\begin{theorem}[Verifier call bound]
\label{thm:calls}
Under the same conditions as Theorem~\ref{thm:main}, with
\[
  N_{\mathcal{T}} \;=\; \frac{K(K^D-1)}{K-1} \;=\; K + K^2 + \cdots + K^D
\]
non-root nodes and threshold schedule $\tau_t = \gamma/\sqrt{t}$:

\begin{enumerate}[(a)]
  \item \textbf{(Single-verification ceiling)}
        $\mathbb{E}[\mathcal{B}(\pi)] \leq N_{\mathcal{T}} = \dfrac{K(K^D-1)}{K-1}$.

  \item \textbf{(Threshold-triggered refinement)}
        Define $T_{\mathrm{free}} = \lfloor(\gamma/(LD))^2\rfloor$.
        For $t \leq T_{\mathrm{free}}$, $\tau_t \geq LD$ dominates every possible
        depth gap $\delta_t \leq LD$, so \emph{threshold-triggered} calls
        (on re-visited nodes, $s_{\mathrm{anc}} \neq \emptyset$) cannot occur.
        Verifier calls during $t \leq T_{\mathrm{free}}$ are exclusively
        first-visit calls ($s_{\mathrm{anc}} = \emptyset$), already bounded by
        Part~(a).  The number of threshold-triggered calls is at most
        $\max(0,\, T - T_{\mathrm{free}})$.
        $T_{\mathrm{free}}$ grows with the signal-to-noise ratio $\gamma/(LD)$:
        higher $\gamma$ or smaller $L$ extends the first-visit-only phase.

  \item \textbf{(Savings fraction)}
        For any $T > N_{\mathcal{T}}$, the fraction of steps requiring verification satisfies
        \[
          \frac{\mathbb{E}[\mathcal{B}(\pi)]}{T} \;\leq\; \frac{N_{\mathcal{T}}}{T} \;\to\; 0
          \quad\text{as } T \to \infty \text{ with } K, D \text{ fixed}.
        \]
\end{enumerate}
\end{theorem}

\begin{proof}[Proof sketch]
Part (a): DVA-MCTS stores $\hat{V}(s)$ on first verification and reuses it on all
subsequent visits, so any node is verified at most once. Hence $\mathcal{B}(\pi) \leq N_{\mathcal{T}}$.

Part (b): A threshold-triggered call requires both $s_{\mathrm{anc}} \neq \emptyset$ (re-visit)
and $\delta_t > \tau_t$.  Since $\delta_t \leq LD$ always, $\tau_t \geq LD$
(i.e.\ $t \leq T_{\mathrm{free}}$) prevents any threshold-triggered call.
Calls at $t \leq T_{\mathrm{free}}$ are exclusively first-visit calls
($s_{\mathrm{anc}} = \emptyset$), already counted in Part~(a)'s $N_{\mathcal{T}}$ ceiling.

Part (c): Immediate from (a).
The full proof is in Appendix~\ref{app:calls}.
\end{proof}

\begin{remark}[Interpreting the two bounds]
\label{rem:regime_calls}
Part (a) is driven by tree size; part (b) is driven by the algorithm's parameters
$\gamma$ and $L$. For the theoretically optimal $\gamma^* = 4\sigma_{\max}\sqrt{2\log K}$
(Theorem~\ref{thm:main}), $T_{\mathrm{free}} = 32\sigma_{\max}^2\log K / (L^2D^2)$:
a constant that grows with signal-to-noise $\sigma_{\max}/L$ and shrinks with tree
depth $D$. Part (b) thus shows that a \emph{better signal-to-noise ratio or shallower
tree} yields a longer \emph{first-visit-only} initial phase---one in which
threshold-triggered re-visit calls cannot occur---even within the exploration regime.
Three distinct regimes emerge empirically (see also Section~\ref{sec:regime}):
\emph{exploration} ($B \ll N_{\mathcal{T}}$, savings ${<}5\%$);
\emph{transition} ($B \lesssim N_{\mathcal{T}}$, proxy mechanism active,
savings $23\%$--$73\%$ for $B \in [800, 3200]$);
and \emph{exploitation} ($B \gg N_{\mathcal{T}}$, single-verification ceiling active,
savings $\to N_{\mathcal{T}}/B \to 0$: $93.7\%$ at $B{=}16{,}384$,
and $94.6\%$ at $B{=}400$ for the small-tree setting $K{=}2$, $D{=}4$).
In the exploitation regime, both parts (a) and (b) of Theorem~\ref{thm:calls} agree
and the savings fraction $N_{\mathcal{T}}/T \to 0$
(Section~\ref{sec:exploitation_results}).
\end{remark}

\subsection{Lower Bound}

\begin{theorem}[Minimax lower bound]
\label{thm:lower}
For any online verifier allocation policy $\pi$, there exists a problem instance
(a tree with branching factor $K$ and verifier noise $\sigma > 0$) such that
\[
  \mathbb{E}[R(T)] \;\geq\; \frac{\sigma}{8}\sqrt{T}.
\]
\end{theorem}

\begin{proof}[Proof sketch]
We construct a hard instance via a two-node tree ($K=2$). One node has true value
$V^* = 1/2 + \Delta$ and the other $V = 1/2 - \Delta$. The optimal policy must
distinguish these two nodes. By Le Cam's method \citep{lecam1973convergence}, any
test that distinguishes the nodes with probability $\geq 3/4$ requires at least
$\Omega(1/\Delta^2)$ verifier calls. Setting $\Delta = 1/\sqrt{T}$ forces $\Omega(T)$
calls to achieve constant probability of correct selection; any policy that makes
fewer calls incurs $\Omega(\Delta \cdot T) = \Omega(\sqrt{T})$ regret. The full
argument via Pinsker's inequality and Fano's method is in Appendix~\ref{app:lower}.
\end{proof}

\begin{corollary}[Rate-optimality]
\label{cor:optimal}
DVA-MCTS achieves $O(\sqrt{T \log K})$ regret (Theorem~\ref{thm:main}) while the
minimax lower bound (Theorem~\ref{thm:lower}) gives $\Omega(\sqrt{T})$, establishing
rate-optimality up to the $\sqrt{\log K}$ factor.
The constant $C(K,D)$ is conservative by a factor $\Phi(K,D)$ relative to the tight
$D{=}1$ value; this is a proof artifact (Remark~\ref{rem:constant_tightness}),
not a property of the algorithm.
\end{corollary}

\subsection{Finite-Sample Behaviour and the Exploration Regime}

\begin{remark}[Three budget regimes]
\label{rem:exploration}
Theorem~\ref{thm:main} is an asymptotic statement as $T \to \infty$.
At finite $T$, the algorithm passes through three distinct regimes determined by
$N_{\mathcal{T}} = K(K^D-1)/(K-1)$ (Section~\ref{sec:regime}):

\begin{enumerate}[(i)]
  \item \textbf{Exploration} ($T \ll N_{\mathcal{T}}$): Most visited nodes are new.
    DVA calls the verifier on every first visit ($s_{\text{anc}} = \emptyset$),
    call rate $\approx 1$, savings ${<}5\%$.  The log-log regret slope exceeds $0.5$
    due to $\Theta(T)$ first-visit overhead; this bias vanishes as
    $T/N_{\mathcal{T}} \to \infty$.

  \item \textbf{Transition} ($T \lesssim N_{\mathcal{T}}$): The tree is partially
    covered.  DVA's Lipschitz proxy substitutes for verifier calls at first-visit
    nodes whose parent score is within threshold $\tau_t$.  Savings grow from
    ${\sim}5\%$ to ${\sim}73\%$ across $B \in [400, 3200]$ (proxy mechanism).

  \item \textbf{Exploitation} ($T \gg N_{\mathcal{T}}$): Most nodes have stored
    observations.  The single-verification ceiling
    (Theorem~\ref{thm:calls}(a)) dominates: call count $\leq N_{\mathcal{T}}$,
    savings fraction $N_{\mathcal{T}}/T \to 0$, and the $O(\sqrt{T\log K})$
    rate emerges cleanly.
\end{enumerate}

In the main experiments ($K{=}2$, $D{=}12$, $N_{\mathcal{T}}{=}8{,}190$),
budgets $B \leq 400$ are firmly in the exploration regime (explains slope $0.635$
at $T{=}400$); $B \in [800, 3200]$ are in the transition regime; and $B{=}16{,}384$
first enters the exploitation regime.  The small-tree setting ($K{=}2$, $D{=}4$,
$N_{\mathcal{T}}{=}30$) makes the exploitation plateau accessible at $B{=}100$
(Section~\ref{sec:small_tree_results}).
\end{remark}

\begin{remark}[Sensitivity to Lipschitz constant estimation]
\label{rem:lhat}
The algorithm takes $L$ as input. In practice $L$ must be estimated; the natural
estimator $\hat{L} = \mathbb{E}[|V(s){-}V(s')|] / |d{-}d'|$ underestimates the
supremum because it measures the \emph{mean} rate of change. Plugging $\hat{L} < L$
into DVA-MCTS reduces the threshold $\delta_t = \hat{L} \cdot |d{-}d_{\text{anc}}|$,
causing the algorithm to trigger fewer verifier calls.  In the \emph{exploitation
regime} ($T \gg N_{\mathcal{T}}$), this means the algorithm skips some high-variance
transitions that should be re-verified, degrading the regret bound by a factor of
$(L / \hat{L})$ on the estimation-error term.

Two estimation methods are studied empirically (Section~\ref{sec:lhat}).
The \emph{global pair estimator} $\hat{L}_{\text{pair}} = \mathbb{E}[|V(s){-}V(s')|/|d{-}d'|]$
yields $\hat{L}_{\text{pair}}{=}0.091$ ($4\times$ underestimate) due to multi-step
cancellation bias.  The \emph{sequential step estimator}
$\hat{L}_{\text{seq}} = \mathbb{E}[|V(s){-}V(s_{\text{prev}})|]$ yields
$\hat{L}_{\text{seq}}{=}0.139$ ($2.5\times$ underestimate).
In the exploration regime ($T{=}400$, Table~\ref{tab:lhat}), differences are
statistically non-significant (Wilcoxon $p{=}0.35$): fewer verifier calls reduce
noise accumulation and can lower cumulative regret.  The exploitation regime,
where threshold quality determines re-verification decisions, is expected to show
the theoretically predicted degradation.
An adaptive estimator tracking the running maximum of $|V(s){-}V(s')| / |d{-}d'|$
gives a consistent online estimate of $L$ that converges to $\hat{L}{=}0.431$
(a conservative $23\%$ over-estimate), achieves the best empirical regret in the
exploration regime, and is the recommended default
(Section~\ref{sec:lhat}).
\end{remark}

\subsection{Robustness to Assumption Violations}

\begin{remark}[Robustness to Lipschitz violations]
\label{rem:robustness}
Suppose the Lipschitz condition (Assumption~\ref{ass:lipschitz}) holds everywhere
except on a set $\mathcal{E}$ of ``exception'' transitions with $|\mathcal{E}| = m$.
Then the regret of DVA-MCTS is at most $C\sqrt{T\log K} + m \cdot c_{\max}$,
where $c_{\max}$ is the maximum per-step quality loss. For $m = O(\sqrt{T})$,
this preserves the $O(\sqrt{T \log K})$ rate.
\end{remark}

\subsection{Connection to Compute-Optimal Scaling}

\citet{snell2024scaling} identifies the compute-optimal strategy as a function of
problem difficulty. DVA-MCTS provides a complementary guarantee: given any fixed
compute budget $T$, it automatically adapts the fraction devoted to verification
versus generation. For easy problems (low score variance), $\delta_t$ rarely
exceeds $\tau_t$, so very few verifier calls are made. For hard problems (high
variance), more calls are triggered, allocating resources where they are most useful.
This is consistent with the difficulty-adaptive compute allocation of
\citet{wu2024inference} but with formal guarantees.
